\documentclass[11pt]{article}
% Shared print layout. The web's provisional design and MathJax are independent.
\usepackage[utf8]{inputenc}
\usepackage[T1]{fontenc}
\usepackage[spanish,es-nodecimaldot]{babel}
\usepackage{lmodern}
\usepackage{amsmath,amssymb}
\usepackage{graphicx}
\usepackage{geometry}
\geometry{a4paper,left=20mm,right=20mm,top=25mm,bottom=22mm,headheight=16pt,headsep=8mm}
\usepackage{fancyhdr,lastpage}
\usepackage{enumitem}
\usepackage{xcolor}
\usepackage{pgfplots}
\pgfplotsset{compat=1.18}
\usepgfplotslibrary{fillbetween}
\usepackage[hidelinks,pdfencoding=auto]{hyperref}
\hypersetup{pdflang=es-ES,pdfauthor={},pdfsubject={Matemáticas Aplicadas a las Ciencias Sociales II}}
\pagestyle{fancy}
\fancyhf{}
\fancyhead[L]{\small Matemáticas Aplicadas a las Ciencias Sociales II}
\fancyhead[R]{\small 2.\textsuperscript{o} Bachillerato}
\fancyfoot[L]{\footnotesize Inferencia II | Edición revisada 1.0}
\fancyfoot[R]{\footnotesize Página \thepage\ de \pageref{LastPage}}
\renewcommand{\headrulewidth}{0.4pt}
\setlength{\parindent}{0pt}
\setlength{\parskip}{0.5em}
\setlist{itemsep=0.3em,topsep=0.3em}
\newcommand{\apartado}[1]{\subsection*{#1}}
\pgfplotsset{
  revision/.style={
    width=0.94\linewidth,height=7.5cm,
    axis lines=middle,grid=major,grid style={gray!18},
    tick label style={font=\small},label style={font=\small},
    xlabel={$x$},ylabel={$y$},samples=160,
    legend style={at={(0.5,-0.16)},anchor=north,draw=none,font=\small},
    clip=true
  }
}

\hypersetup{pdftitle={Inferencia II - Examen revisado}}
\begin{document}
\begin{minipage}{0.17\linewidth}\includegraphics[width=0.8\linewidth]{logo-al-andalus.png}\end{minipage}\hfill
\begin{minipage}{0.79\linewidth}{\Large\bfseries Inferencia II}\\[0.4em]
Duración: \textbf{55 minutos}. Puntuación: \textbf{10 puntos}.\end{minipage}

Apellidos y nombre: \hrulefill

Fecha: \hrulefill

Responde razonadamente y numera tus páginas. Se permite calculadora. Cada ejercicio incluye 0,5 puntos de expresión de procesos y métodos; no se añaden a su puntuación máxima. Se conservan los tiempos orientativos originales de 20, 20 y 10 minutos, dentro de los 55 minutos totales.


\section*{Ejercicio 1 \hfill 4 puntos} \textbf{(20')} Se desea estimar el porcentaje de jóvenes que utilizan una determinada red social. Para ello se escoge una muestra aleatoria simple de 500 jóvenes y de ellos 410 afirman utilizarla.

\begin{enumerate}[label=(\alph*)]
  \item (2,5 puntos)  Calcule el intervalo de confianza para la proporción de jóvenes que usa esa red social con un nivel de confianza del 95\%.

  

  \item (1,5 puntos)  Manteniendo la proporción muestral, determine el tamaño mínimo de la muestra necesario para que, con un nivel de confianza del 97\%, el error máximo que se cometa al estimar la proporción de esa población sea inferior a 0.04.

  

\end{enumerate}


\section*{Ejercicio 2 \hfill 4 puntos} \textbf{(20')} El precio de ciertos electrodomésticos se puede considerar como una variable aleatoria con distribución normal y desviación típica conocida de 100 euros. Se toman los precios (en euros) de una muestra de 9 electrodomésticos:

\[
255,\ 85,\ 120,\ 290,\ 80,\ 90,\ 275,\ 290,\ 135.
\]

\begin{enumerate}[label=(\alph*)]

 \item (2,5 puntos)  Construye un intervalo de confianza al 95\% para la media poblacional.



  \item (1,5 puntos)  Halla el tamaño mínimo que debe tener la muestra para que, con un nivel de confianza del 99\%, el error de estimación del precio medio no supere los 50 euros.

  

\end{enumerate}


\clearpage
\section*{Ejercicio 3 \hfill 2 puntos} \textbf{(10')} Una universidad desea realizar un estudio sobre los hábitos culturales del alumnado, dividiéndolos en cuatro estratos según la facultad a la que pertenecen: Ciencias Sociales, Ciencias de la Salud, Ingeniería y Humanidades, con tamaños de 4800, 2200, 3000 y 2000 estudiantes, respectivamente. Se realiza un muestreo estratificado con afijación proporcional y se han seleccionado 60 estudiantes de la Facultad de Ingeniería.

\begin{enumerate}[label=(\alph*)]
  \item (1 punto)  Determina el tamaño total de la muestra y cuántos estudiantes deben seleccionarse de cada estrato.
  
  

  \item (1 punto)  Explica la diferencia entre muestreo probabilístico y no probabilístico. Pon un ejemplo de cada tipo aplicado al estudio de hábitos culturales en la universidad, y reflexiona: \textquestiondown{}cuál crees que ofrece resultados más precisos? Justifica tu respuesta.

  
\end{enumerate}





\end{document}
